% GENERATED FILE. Edit canonical lessons, then run npm run generate:books.
% canonical-source-hash: fnv1a64:d6cd99935364db30
% canonical-lessons: RU-C361-predpriyatie, RU-C361-pribyl, RU-C361-ubytok, RU-C361-dokhod, RU-C361-sdelka

\chapter{Profit and loss}
\label{ch:ru-profit-and-loss}
\hlchaptermodality{361}

\begin{chapteropening}
\textbf{By the end of this chapter:} \emph{I can talk about a business, its profit and its losses, someone’s income, and a deal.}

Complete the last lesson of chapter 361: I can talk about a business, its profit and its losses, someone’s income, and a deal.
\end{chapteropening}

\section[\texorpdfstring{\ru{предприятие} (predpriyátiye)}{предприятие (predpriyátiye)}]{\ru{предприятие} (predpriyátiye) — an enterprise, a business; a plant}

\label{lesson:RU-C361-predpriyatie}



\begin{quote}
Before the new one: say the Russian for a shift, then the Russian for to report.
\end{quote}

\subsection*{You'll want to know: \ru{предприятие}}
\textbf{\ru{предприятие}} (\emph{predpriyátiye}) — \textquotedblleft{}an enterprise, a business; a plant\textquotedblright{}.

\begin{quote}
\ru{Отец} \ru{работает} \ru{на} \ru{большом} \ru{предприятии}.
\end{quote}

\begin{quote}
\emph{Otéts rabótayet na bolshóm predpriyátii.}
\end{quote}

\begin{quote}
My father works at a large plant.
\end{quote}

Neuter. \emph{\ru{На} \ru{предприятии}} — at a plant. \emph{\ru{Малое} \ru{предприятие}} — a small business.

\begin{cousinweb}
\textbf{\ru{предприятие}} is built on \emph{\ru{предпринять}} \textquotedblleft{}to undertake\textquotedblright{}, from \emph{\ru{пред}-} \textquotedblleft{}before\textquotedblright{} and the root of \emph{\ru{принять}} \textquotedblleft{}to take\textquotedblright{}, from the old verb \textquotedblleft{}to take\textquotedblright{} seen in \emph{\ru{снимать}}: an undertaking.
\end{cousinweb}

\begin{grammarlens}[title={What you've built}]
One more word for telling, explaining and describing at B1.
\end{grammarlens}

\subsection*{Your turn}
\begin{itemize}
\raggedright
  \item \emph{Say it:} \emph{predpriyátiye}
  \item \emph{Say it:} \emph{predpriyátiye}, once more
  \item \emph{Say it:} a sentence of your own with \emph{predpriyátiye}
  \item \emph{Recall:} say the Russian for a shift, then the Russian for to report, then say \emph{predpriyátiye} again
\end{itemize}

\begin{tcolorbox}[breakable,colback=teal!4,colframe=teal!35!black,title={Before you move on}]
What does \ru{предприятие} mean? (\textquotedblleft{}An enterprise, a business; a plant\textquotedblright{}.) Say it once more.
\end{tcolorbox}
\section[\texorpdfstring{\ru{прибыль} (príbyl)}{прибыль (príbyl)}]{\ru{прибыль} (príbyl) — profit}

\label{lesson:RU-C361-pribyl}



\begin{quote}
Before the new one: say the Russian for to report, then the Russian for an enterprise.
\end{quote}

\subsection*{You'll want to know: \ru{прибыль}}
\textbf{\ru{прибыль}} (\emph{príbyl}) — \textquotedblleft{}profit\textquotedblright{}.

\begin{quote}
\ru{В} \ru{этом} \ru{году} \ru{фирма} \ru{получила} \ru{большую} \ru{прибыль}.
\end{quote}

\begin{quote}
\emph{V étom godú fírma poluchíla bolshúyu príbyl.}
\end{quote}

\begin{quote}
This year the firm made a large profit.
\end{quote}

Feminine. \emph{\ru{Приносить} \ru{прибыль}} — to make a profit.

\begin{cousinweb}
\textbf{\ru{прибыль}} is built on \emph{\ru{прибыть}} \textquotedblleft{}to arrive, to increase\textquotedblright{}, from \emph{\ru{при}-} and \emph{\ru{быть}} \textquotedblleft{}to be\textquotedblright{}, from Proto-Slavic \emph{*byti}: what comes to be in addition.
\end{cousinweb}

\begin{grammarlens}[title={What you've built}]
One more word for telling, explaining and describing at B1.
\end{grammarlens}

\subsection*{Your turn}
\begin{itemize}
\raggedright
  \item \emph{Say it:} \emph{príbyl}
  \item \emph{Say it:} \emph{príbyl}, once more
  \item \emph{Say it:} a sentence of your own with \emph{príbyl}
  \item \emph{Recall:} say the Russian for to report, then the Russian for an enterprise, then say \emph{príbyl} again
\end{itemize}

\begin{tcolorbox}[breakable,colback=teal!4,colframe=teal!35!black,title={Before you move on}]
What does \ru{прибыль} mean? (\textquotedblleft{}Profit\textquotedblright{}.) Say it once more.
\end{tcolorbox}
\section[\texorpdfstring{\ru{убыток} (ubýtok)}{убыток (ubýtok)}]{\ru{убыток} (ubýtok) — a loss (of money)}

\label{lesson:RU-C361-ubytok}



\begin{quote}
Before the new one: say the Russian for an enterprise, then the Russian for profit.
\end{quote}

\subsection*{You'll want to know: \ru{убыток}}
\textbf{\ru{убыток}} (\emph{ubýtok}) — \textquotedblleft{}a loss (of money)\textquotedblright{}.

\begin{quote}
\ru{Магазин} \ru{работает} \ru{в} \ru{убыток}.
\end{quote}

\begin{quote}
\emph{Magazín rabótayet v ubýtok.}
\end{quote}

\begin{quote}
The shop is running at a loss.
\end{quote}

Masculine; genitive \emph{\ru{убытка}}. Usually plural: \emph{\ru{большие} \ru{убытки}}.

\begin{cousinweb}
\textbf{\ru{убыток}} is built on \emph{\ru{убыть}} \textquotedblleft{}to decrease\textquotedblright{}, from \emph{\ru{у}-} \textquotedblleft{}away\textquotedblright{} and \emph{\ru{быть}} \textquotedblleft{}to be\textquotedblright{}, from Proto-Slavic \emph{*byti}: what goes away.
\end{cousinweb}

\begin{grammarlens}[title={What you've built}]
One more word for telling, explaining and describing at B1.
\end{grammarlens}

\subsection*{Your turn}
\begin{itemize}
\raggedright
  \item \emph{Say it:} \emph{ubýtok}
  \item \emph{Say it:} \emph{ubýtok}, once more
  \item \emph{Say it:} a sentence of your own with \emph{ubýtok}
  \item \emph{Recall:} say the Russian for an enterprise, then the Russian for profit, then say \emph{ubýtok} again
\end{itemize}

\begin{tcolorbox}[breakable,colback=teal!4,colframe=teal!35!black,title={Before you move on}]
What does \ru{убыток} mean? (\textquotedblleft{}A loss (of money)\textquotedblright{}.) Say it once more.
\end{tcolorbox}
\section[\texorpdfstring{\ru{доход} (dokhód)}{доход (dokhód)}]{\ru{доход} (dokhód) — income}

\label{lesson:RU-C361-dokhod}



\begin{quote}
Before the new one: say the Russian for profit, then the Russian for a loss.
\end{quote}

\subsection*{You'll want to know: \ru{доход}}
\textbf{\ru{доход}} (\emph{dokhód}) — \textquotedblleft{}income\textquotedblright{}.

\begin{quote}
\ru{У} \ru{них} \ru{небольшой} \ru{доход}.
\end{quote}

\begin{quote}
\emph{U nikh nebolshóy dokhód.}
\end{quote}

\begin{quote}
They have a small income.
\end{quote}

Masculine. \emph{\ru{Доходы} \ru{и} \ru{расходы}} — income and expenditure.

\begin{cousinweb}
\textbf{\ru{доход}} is built from \emph{\ru{до}-} \textquotedblleft{}up to\textquotedblright{} and \emph{\ru{ход}} \textquotedblleft{}a going\textquotedblright{}, from \emph{\ru{ходить}} \textquotedblleft{}to go\textquotedblright{}, Proto-Slavic \emph{*xoditi}: what comes up to you.
\end{cousinweb}

\begin{grammarlens}[title={What you've built}]
One more word for telling, explaining and describing at B1.
\end{grammarlens}

\subsection*{Your turn}
\begin{itemize}
\raggedright
  \item \emph{Say it:} \emph{dokhód}
  \item \emph{Say it:} \emph{dokhód}, once more
  \item \emph{Say it:} a sentence of your own with \emph{dokhód}
  \item \emph{Recall:} say the Russian for profit, then the Russian for a loss, then say \emph{dokhód} again
\end{itemize}

\begin{tcolorbox}[breakable,colback=teal!4,colframe=teal!35!black,title={Before you move on}]
What does \ru{доход} mean? (\textquotedblleft{}Income\textquotedblright{}.) Say it once more.
\end{tcolorbox}
\section[\texorpdfstring{\ru{сделка} (sdélka)}{сделка (sdélka)}]{\ru{сделка} (sdélka) — a deal, a transaction}

\label{lesson:RU-C361-sdelka}



\begin{quote}
Before the new one: say the Russian for a loss, then the Russian for income.
\end{quote}

\subsection*{You'll want to know: \ru{сделка}}
\textbf{\ru{сделка}} (\emph{sdélka}) — \textquotedblleft{}a deal, a transaction\textquotedblright{}.

\begin{quote}
\ru{Компании} \ru{заключили} \ru{сделку}.
\end{quote}

\begin{quote}
\emph{Kompánii zaklyuchíli sdélku.}
\end{quote}

\begin{quote}
The companies struck a deal.
\end{quote}

Feminine. \emph{\ru{Заключить} \ru{сделку}} — to strike a deal.

\begin{cousinweb}
\textbf{\ru{сделка}} is built on \emph{\ru{сделать}} \textquotedblleft{}to do, to make\textquotedblright{}, from \emph{\ru{делать}}, inherited from Proto-Slavic \emph{*dělati}: a deal is a thing done between two sides.
\end{cousinweb}

\begin{grammarlens}[title={What you've built}]
One more word for telling, explaining and describing at B1.
\end{grammarlens}

\subsection*{Your turn}
\begin{itemize}
\raggedright
  \item \emph{Say it:} \emph{sdélka}
  \item \emph{Say it:} \emph{sdélka}, once more
  \item \emph{Say it:} a sentence of your own with \emph{sdélka}
  \item \emph{Recall:} say the Russian for a loss, then the Russian for income, then say \emph{sdélka} again
\end{itemize}

\begin{tcolorbox}[breakable,colback=teal!4,colframe=teal!35!black,title={Before you move on}]
What does \ru{сделка} mean? (\textquotedblleft{}A deal, a transaction\textquotedblright{}.) Say it once more.
\end{tcolorbox}
